% ============================================================
% ch17.tex —— 第 17 章 堆积的艺术：发明积分
% ============================================================
\chapter{堆积的艺术：发明积分}

三角形、长方形、梯形都有面积公式——因为边是直的。可抛物线 $y = x^2$ 下方那块"底边直、顶上弯"的月牙地，面积是多少？阿基米德为之绞尽脑汁（他用的是双重反证法，写了几十页）；你现在要用刘徽的刀（第 14 章）重新挣一遍这块地——顺便把圆面积公式也重新挣一遍。

\begin{kaiwei}
曲线 $y = x^2$ 下方、$x$ 从 $0$ 到 $1$ 的区域，面积是多少？（先猜：它比底 $1$ 高 $1$ 的三角形面积 $\tfrac12$ 大还是小？）
\end{kaiwei}

直觉核对，把范围先圈死：在 $0 < x < 1$ 上恒有 $x^2 < x$（一个真分数的平方比它自己小），所以曲边 $y = x^2$ 整个位于直线 $y = x$ \textbf{下方}——月牙地完全装在"底 $1$ 高 $1$ 的三角形"里。于是
\[
0 \;<\; \text{月牙地} \;<\; \frac12 .
\]
圈定嫌疑范围，接下来精确破案。

\section{笨办法先行：切条}

思路朴素到近乎无赖：\textbf{不会算弯的，就切成直的}。把 $[0,1]$ 切成 $n$ 条等宽竖条，每条宽 $\tfrac1n$；每条上方的曲面顶棚换成平顶——取\textbf{右端点}处的高度。第 $k$ 条（$k = 1, \dots, n$）的右端点横坐标是 $\tfrac kn$，高 $\left(\tfrac kn\right)^2$：

\begin{tikzpicture}[scale=3.2]
  \draw[->] (-0.15,0) -- (1.3,0) node[right] {$x$};
  \draw[->] (0,-0.15) -- (0,1.25) node[above] {$y$};
  \foreach \k in {1,...,4} {
    \draw[fill=cheng!18] ({(\k-1)/4},0) rectangle ({\k/4},{\k*\k/16});
  }
  \draw[thick,domain=0:1.05,smooth] plot (\x,{\x*\x});
  \draw[dashed] (0.25,0.0625) -- (0.25,0.25);
  \draw[dashed] (0.5,0.25) -- (0.5,0.5);
  \draw[dashed] (0.75,0.5625) -- (0.75,0.75);
  \node at (1.18,1.08) {\scriptsize $y=x^2$};
  \node at (0.5,-0.32) {\scriptsize 切四条的阶梯（右端点取平顶，全部高估）};
\end{tikzpicture}

切 $n = 4$ 条先试手感：四条高度依次为 $\left(\tfrac14\right)^2, \left(\tfrac24\right)^2, \left(\tfrac34\right)^2, 1$，阶梯面积
\[
S_4 = \frac14\left(\frac{1}{16} + \frac{4}{16} + \frac{9}{16} + \frac{16}{16}\right) = \frac{30}{64} = 0.469 .
\]
这个 $0.469$ 是\textbf{高估}：曲线上升，右端点是每条内的最高点，平顶不低于曲顶，所以阶梯盖住月牙地——月牙地 $< 0.469 < 0.5$，与刚才圈定的范围吻合。切细一点：$n = 10$，$S_{10} = 0.385$；$n = 100$，$S_{100} = 0.338$——数字一路下滑，奔着某个数去。

一般地，切 $n$ 条（右端点平顶）：
\[
S_n = \frac1n\left[\left(\frac1n\right)^2 + \left(\frac2n\right)^2 + \cdots + \left(\frac nn\right)^2\right]
= \frac{1^2 + 2^2 + \cdots + n^2}{n^3} .
\]
分子是平方和——它有闭式公式（初中可验证）：
\[
1^2 + 2^2 + \cdots + n^2 = \frac{n(n+1)(2n+1)}{6} \qquad
\text{（}n = 3\text{ 验证：}1+4+9 = 14 = \tfrac{3\cdot4\cdot7}{6}\text{ \checkmark）} .
\]
代入整理：
\[
S_n = \frac{n(n+1)(2n+1)}{6n^3} = \frac{2n^3 + 3n^2 + n}{6n^3} = \frac13 + \frac{1}{2n} + \frac{1}{6n^2} .
\]
（上下同除 $n^3$——第 14 章双扳手之一。）让 $n$ 进场比赛：后两项压成零，
\[
\boxed{\;\lim_{n\to\infty} S_n = \frac13\;}
\]
切一万条、切一亿条，答案纹丝不动地停在 $\tfrac13$——\textbf{月牙地的面积是 $\tfrac13$}。

\textbf{两侧夹击，防止侥幸}。右端点平顶处处不低于曲线（$x^2$ 上升，右端是条内最高点），所以 $S_n$ 是\textbf{高估}。换\textbf{左端点}平顶（处处低估，第 $k$ 条取高 $\left(\tfrac{k-1}n\right)^2$）：
\[
S_n^{-} = \frac{(0)^2 + 1^2 + \cdots + (n-1)^2}{n^3} = \frac{(n-1)n(2n-1)}{6n^3} = \frac13 - \frac{1}{2n} + \frac{1}{6n^2} .
\]
（同款公式，$n$ 换 $n-1$。）两串阶梯一上一下：
\[
S_n^{-} < \text{真面积} < S_n, \qquad \text{而两串都奔 } \tfrac13 .
\]
\textbf{夹逼定理}（第 14 章烧脑时刻你证过）出庭：真面积无处可逃，恰为 $\tfrac13$。这与刘徽"夹 $3.141024 < \pi < 3.142704$"是同一招——\textbf{上阶梯与下阶梯是现代版的"割圆内外多边形"}。

\begin{faming}
\textbf{土名}：堆积、切条求和 $\to$ \textbf{正式名}：\textbf{定积分}（definite integral），记作
\[
\int_0^1 x^2\,dx = \frac13 .
\]
拉长的"S"（$\int$，莱布尼茨的手笔）是拉丁文 \textit{summa}（和）的首字母——它就是"$S_n$ 的无穷延续"；$"dx"$ 是"每条的宽"（$h\to0$ 的化身）；下上标 $0, 1$ 是左右边界。整串读作："$x^2$ 从 $0$ 积到 $1$，结果 $\tfrac13$。"
\end{faming}

\section{堆积的另一副面孔：从速度表算里程}

面积只是开胃菜。堆积真正的杀手级应用是\textbf{把无穷多个瞬间的变化累成总变化}。

第 15 章走了一遍"路程 $\to$ 瞬间速度"（求导）。现在\textbf{倒着走}：只给你速度函数，算总路程。设速度 $v(t) = 2t$（正是 $s = t^2$ 的导数），求 $0$ 到 $3$ 秒走了多远。

匀速的话"路程 $=$ 速度 $\times$ 时间"一条公式收工；变速不行——那就切时间！把 $[0,3]$ 切成 $n$ 段，每段长 $\tfrac3n$。第 $k$ 段内速度近似不变（取右端点值 $v\left(\tfrac{3k}n\right) = \tfrac{6k}n$），该段路程 $\approx \dfrac{6k}{n}\cdot\dfrac3n$。总路程：
\[
R_n = \sum_{k=1}^n \frac{6k}{n}\cdot\frac3n = \frac{18}{n^2}(1 + 2 + \cdots + n) = \frac{18}{n^2}\cdot\frac{n(n+1)}{2} = 9 + \frac{9}{n} .
\]
（等差数列求和，初中存货。）比赛：$R_n \to 9$。\textbf{$0$ 到 $3$ 秒走了 $9$ 米}。

验算：如果路程函数真的是 $s = t^2$，那 $s(3) - s(0) = 9 - 0 = 9$。\textbf{严丝合缝}。

这不是巧合，这是结构：\textbf{堆积（速度 $\to$ 路程）与求导（路程 $\to$ 速度）互为逆操作}。一句话的预感，完整的定理在第 18 章。先记体感：
\[
\int_a^b (\text{速度})\,dt = (\text{路程})\big|_a^b = s(b) - s(a), \qquad
\frac{d}{dt}(\text{路程}) = \text{速度} .
\]
（$\big|$ 是"代入相减"的速记，第 18 章正式颁发。）

\section{重挣一遍圆面积：切在圆心}

刘徽从\textbf{圆周}切多边形（内接逼近），阿基米德同款。换个刀法——\textbf{从圆心切扇形}，看一场漂亮的堆积表演。

把半径 $r$ 的圆切成 $n$ 个细扇形（像切披萨）。每个扇形近似看作"底为弧、高为 $r$"的细长三角——因为顶角极小，两条腰几乎平行，弧底几乎贴直。再把所有扇形\textbf{一上一下交错摆放}：

\begin{tikzpicture}[scale=0.9]
  \foreach \i in {0,...,7} {
    \pgfmathsetmacro\x{0.42*\i}
    \pgfmathsetmacro\y{mod(\i,2)==0 ? 0 : 0.5}
    \draw[fill=shenlan!15] (\x,\y) -- (\x+0.42,\y) -- (\x+0.21,\y+0.55) -- cycle;
  }
  \draw[<->] (0,-0.4) -- (3.4,-0.4) node[midway,below] {\scriptsize 底 $\approx \pi r$（半个周长）};
  \draw[<->] (3.75,0) -- (3.75,1.05) node[midway,right] {\scriptsize 高 $=r$};
\end{tikzpicture}

锯齿边互相咬合，拼出的图形越来越像一个标准\textbf{三角形}：底是半个周长 $\pi r$（$n$ 个弧底轮流朝上朝下，摊成两条近似直的长边），高是 $r$。扇形切得越细，锯齿越平，拼图面积的落脚点：
\[
\text{圆面积} = \lim_{n\to\infty}\left(n\times\text{细扇形}\right) = \frac12 \times \pi r \times r = \pi r^2 .
\]
\textbf{$\pi r^2$ 不再是"背下来的公式"，而是堆积的产物。}注意这场证明顺便也解释了 $\pi$ 的另一个身份：$\pi = \dfrac{\text{周长}}{\text{直径}}$ 与 $\pi = \dfrac{\text{面积}}{\text{半径平方}}$ 是同一个数——因为面积本来就是周长"一层层摊平堆积"的结果。（阿基米德《圆的度量》用了类似的双重夹逼，加上外切多边形，把 $\pi$ 夹在 $3\tfrac{10}{71}$ 与 $3\tfrac17$ 之间。）

\section{积分的"单位"直觉}

一个帮你终身不忘的观察：\textbf{积分会换单位}。速度（米/秒）对时间（秒）积分，得路程（米）——"每秒多少"乘"多少秒"，单位相消。曲边高度（米）对底边（米）积分，得面积（平方米）。导数反过来："总量对谁的变化率"，单位是"总量单位 $\div$ 那个轴的单位"。以后见到任何 $\int$，先扫一眼单位，就知道它在物理上"称"的是什么。

\begin{dongshou}
\begin{itemize}
  \yisi 左端点重算 $\displaystyle\int_0^1 x^2\,dx$：写出 $S_n^- = \tfrac13 - \tfrac{1}{2n} + \tfrac{1}{6n^2}$ 的推导（平方和公式代入 $n-1$），并取 $n = 4$ 算出 $S_4^- = 0.219$，与 $S_4 = 0.469$ 一起把真值 $\tfrac13 = 0.333$ 夹在中间 \checkmark。
  \yier 硬算 $\displaystyle\int_0^2 t\,dt$：切 $n$ 条、右端点取高（$\tfrac{2k}n$）、等差求和、取落脚点。答案应为 $2$——先用几何（三角形底 $2$ 高 $2$，面积 $\tfrac12\times2\times2$）验算再动手，两路对账。
  \yier 用堆积法求刹车距离：车速 $v = 10 - t$（米/秒，每秒减 $1$），从踩闸到停（$t = 10$ 秒）。切条、求和、取极限（也可直接几何：直线下三角形面积 $\tfrac12\times10\times10 = 50$ 米）。答案 $50$ 米。
  \yier 用扇形拼图的思路"重挣"一个冷门公式：椭圆面积 $= \pi ab$（$a, b$ 是半长轴半短轴）。提示：把单位圆沿横向均匀拉伸 $a$ 倍、纵向拉伸 $b$ 倍得到椭圆；面积被拉伸 $ab$ 倍；圆面积 $\pi$ 拉成 $\pi ab$。（"拉伸改变面积、不改变'圆的积性'"——这是大学"雅可比"思想的初等预告。）
\end{itemize}
\end{dongshou}

\begin{shaonao}
\textbf{硬算 $\displaystyle\int_0^1 t^3\,dt$}。需要立方和公式
\[
1^3 + 2^3 + \cdots + n^3 = \left[\frac{n(n+1)}{2}\right]^2
\]
（它的漂亮之处：立方和恰是"和的平方"。先验证：$1 + 8 = 9 = (1+2)^2$ \checkmark；$1+8+27 = 36 = (1+2+3)^2$ \checkmark。想自己证的话，对 $n$ 归纳，两行。）于是
\[
S_n = \frac{1}{n^4}\cdot\frac{n^2(n+1)^2}{4} = \frac14\left(1 + \frac1n\right)^2 \;\longrightarrow\; \frac14 .
\]
到此为止你已硬算出 $\int_0^1 t^n\,dt = \tfrac1{n+1}$ 的 $n = 1, 2, 3$ 三票（$n=1$：$\tfrac12$；$n=2$：$\tfrac13$；$n=3$：$\tfrac14$）。\textbf{不用平方和公式、立方和公式，一个模式通吃所有 $n$}的方法存在吗？——存在，而且一行：猜原函数 $t^{n+1}/(n+1)$ 求导验证。这就是下一章的后门（牛顿--莱布尼茨公式），你现在站在门缝前。
\end{shaonao}

\begin{chengshi}
"任你怎么切（不均匀切、任意取点）落脚点都相同"是积分合法性的核心（大学叫"黎曼可积"），本章用"左右端点上下夹逼"处理了单调光滑函数，够本书全程使用；一般函数（剧烈振荡、间断）的积分要更细的定义。扇形拼图论证里"锯齿咬合趋于三角形"的严格化同样是极限练习。阿基米德的穷竭法与刘徽的割圆术在严格性上与今天的极限语言等价（差一层记号的皮），两人的年代相差约四百年，方法独立生长——好的数学不止一次被发现。
\end{chengshi}

切条是笨功夫——每算一个积分都要重新切一遍。有没有后门？下一章，微积分最壮丽的定理登场：\textbf{一次发明，终身免切}。
